\documentclass[10pt]{article}
% Surcouche compacte pour les interros
% Socle commun dérivé du framework des khôlles SI
% Version autonome et compatible pdfLaTeX pour le dépôt Interros.
\RequirePackage[T1]{fontenc}
\RequirePackage[utf8]{inputenc}
\RequirePackage[french]{babel}
\RequirePackage{lmodern}
\RequirePackage[margin=1.8cm]{geometry}
\RequirePackage{xcolor}
\RequirePackage{graphicx}
\RequirePackage{amsmath,amssymb}
\RequirePackage{tikz}
\RequirePackage{tabularx}
\RequirePackage{array}
\RequirePackage{enumitem}
\RequirePackage[most]{tcolorbox}
\RequirePackage{fancyhdr}
\RequirePackage{hyperref}
\RequirePackage{titlesec}
\usetikzlibrary{arrows.meta,positioning,calc,fit,chains,decorations.pathmorphing,decorations.markings,shapes,shapes.multipart,patterns,graphs,quotes}
\definecolor{SIblue}{RGB}{31,78,121}
\definecolor{SIlightblue}{RGB}{229,240,250}
\definecolor{SIred}{RGB}{190,0,0}
\definecolor{SIgray}{RGB}{235,235,235}
\hypersetup{colorlinks=true,linkcolor=SIblue,urlcolor=SIblue}
\pagestyle{fancy}
\setlength{\headheight}{14pt}
\fancyhf{}
\lhead{Sciences industrielles pour l'ingénieur}
\rhead{Interros ATS}
\cfoot{\thepage}
\titleformat{\section}{\Large\bfseries\color{SIblue}}{}{0pt}{}
\titleformat{\subsection}{\large\bfseries\color{SIblue}}{}{0pt}{}
\setlist[itemize]{leftmargin=*,itemsep=2pt,topsep=2pt}
\setlist[enumerate]{leftmargin=*,itemsep=4pt,topsep=4pt}
\newtcolorbox{donnees}[1][]{enhanced,breakable,colback=white,colframe=SIblue,title=Données,fonttitle=\bfseries,#1}
\newtcolorbox{attention}[1][]{enhanced,breakable,colback=orange!8,colframe=orange!80!black,title=Attention,fonttitle=\bfseries,#1}
\newtcolorbox{methode}[1][]{enhanced,breakable,colback=green!5,colframe=green!45!black,title=Méthode,fonttitle=\bfseries,#1}
\newcommand{\torseur}[2]{\left\{\mathcal V_{#1}\right\}_{#2}}
% Bibliothèques SI simplifiées, compatibles avec les interros
\providecommand{\og}{``}
\providecommand{\fg}{''}

% Graphe de liaisons
\tikzset{
  glsolide/.style={circle,draw=SIblue,very thick,minimum size=10mm,fill=SIlightblue,font=\bfseries},
  glliaison/.style={draw=SIred,very thick,rounded corners=2pt,font=\small,align=center}
}
\newenvironment{grapheliaisons}[1][]{\begin{center}\begin{tikzpicture}[node distance=25mm,#1]}{\end{tikzpicture}\end{center}}
\newcommand{\GLsolide}[3]{\node[glsolide] (#1) at (#3) {#2};}
\newcommand{\GLliaison}[4][]{\draw[glliaison,#1] (#2) -- node[midway,fill=white,inner sep=2pt]{#4} (#3);}
\newcommand{\GLpivot}[4][]{\GLliaison[#1]{#2}{#3}{Pivot\\$#4$}}
\newcommand{\GLglissiere}[4][]{\GLliaison[#1]{#2}{#3}{Glissière\\$#4$}}
\newcommand{\GLencastrement}[3][]{\GLliaison[#1]{#2}{#3}{Encastrement}}
\newcommand{\GLhelicoidale}[4][]{\GLliaison[#1]{#2}{#3}{Hélicoïdale\\$#4$}}

% Schémas cinématiques simplifiés
\newenvironment{schemacinematique}[1][]{\begin{center}\begin{tikzpicture}[line cap=round,line join=round,#1]}{\end{tikzpicture}\end{center}}
\newcommand{\SCbati}[2]{\coordinate (#1) at (#2);\draw[thick] ($(#1)+(-.65,0)$)--($(#1)+(.65,0)$);\foreach \x in {-.55,-.3,...,.55}{\draw ($(#1)+(\x,0)$)--++(-.18,-.18);}}
\newcommand{\SCpivot}[2]{\node[draw,circle,minimum size=5mm,thick] (#1) at (#2) {};}
\newcommand{\SCglissiere}[2]{\node[draw,rectangle,minimum width=10mm,minimum height=5mm,thick] (#1) at (#2) {};\draw[thick] ($(#1)+(-.35,-.18)$)--($(#1)+(.35,.18)$);}
\newcommand{\SChelicoidale}[2]{\node[draw,circle,minimum size=5mm,thick] (#1) at (#2) {};\draw[thick] ($(#1)+(-.15,0)$) arc (180:-90:.15);}
\newcommand{\SCrelier}[3][]{\draw[very thick,#1] (#2)--(#3);}

% SysML simplifié
\newenvironment{diagrammeSysML}[1][]{\begin{center}\begin{tikzpicture}[node distance=18mm,#1]}{\end{tikzpicture}\end{center}}
\newcommand{\SYSacteur}[4][]{\node[draw,circle,minimum size=3mm,#1] (head#4) at (#2) {};\draw (head#4)--++(0,-.55) coordinate (body#4);\draw ($(body#4)+(0,.35)$)--++(-.28,-.22);\draw ($(body#4)+(0,.35)$)--++(.28,-.22);\draw (body#4)--++(-.25,-.35);\draw (body#4)--++(.25,-.35);\node[below=8mm of head#4] (#4) {#3};}
\newcommand{\SYScas}[4][]{\node[draw,ellipse,align=center,minimum width=28mm,minimum height=10mm,#1] (#4) at (#2) {#3};}
\newcommand{\SYSinclude}[2]{\draw[dashed,->] (#1)--node[above]{\scriptsize\og include\fg}(#2);}
\newcommand{\SYSextend}[2]{\draw[dashed,->] (#1)--node[above]{\scriptsize\og extend\fg}(#2);}

% Tableau de mobilités
\newcommand{\mobilites}[6]{\begin{tabular}{c|c}R&T\\\hline #1&#2\\#3&#4\\#5&#6\end{tabular}}


\geometry{a4paper,left=13mm,right=13mm,top=18mm,bottom=16mm,headheight=28pt,headsep=8mm,footskip=10mm}
\fancyhf{}
\renewcommand{\headrulewidth}{0pt}
\renewcommand{\footrulewidth}{0.6pt}
\fancyfoot[L]{\footnotesize Lycée Robert Doisneau}
\fancyfoot[R]{\footnotesize \thepage}
\newcommand{\InterroHeader}[2]{%
\noindent\begin{tcolorbox}[enhanced,boxrule=.7pt,arc=2mm,colback=SIblue,colframe=SIblue,left=3mm,right=3mm,top=1.5mm,bottom=1.5mm]
\begin{minipage}[c]{.18\linewidth}\color{white}\bfseries\Large SI\end{minipage}%
\begin{minipage}[c]{.60\linewidth}\centering\color{white}\bfseries\large #1\end{minipage}%
\begin{minipage}[c]{.20\linewidth}\raggedleft\color{white}\bfseries Interro\\de cours\end{minipage}
\end{tcolorbox}
\vspace{-1mm}
\noindent\begin{minipage}[c]{.78\linewidth}\begin{tcolorbox}[height=12mm,colback=white,colframe=black,boxrule=1pt,arc=1.5mm,left=2mm,top=2mm]\textbf{Nom :}\end{tcolorbox}\end{minipage}\hfill
\begin{minipage}[c]{.13\linewidth}\begin{tcolorbox}[height=12mm,colback=white,colframe=black,boxrule=1pt,arc=2mm,halign=center,valign=center]\textbf{#2}\end{tcolorbox}\end{minipage}
\vspace{2mm}}
\newcounter{interroquestion}
\newcommand{\question}[2]{\refstepcounter{interroquestion}\par\medskip\noindent\textbf{Question \theinterroquestion :}\quad #1\hfill\textbf{(#2 points)}\par\smallskip}
\newtcolorbox{reponse}[1][35mm]{colback=white,colframe=black,boxrule=.5pt,arc=0mm,height=#1}
\newcommand{\resetquestions}{\setcounter{interroquestion}{0}}
\newcommand{\chaineenergie}{\begin{tikzpicture}[font=\small,>=Latex,node distance=4mm,bloc/.style={draw,minimum width=25mm,minimum height=10mm,align=center}]
\node[bloc](a){\dots};\node[bloc,right=of a](b){MODULER};\node[bloc,right=of b](c){\dots};\node[bloc,right=of c](d){\dots};\node[bloc,fill=yellow!30,right=of d](e){AGIR};
\draw[->,thick](a)--(b);\draw[->,thick](b)--(c);\draw[->,thick](c)--(d);\draw[->,thick](d)--(e);\end{tikzpicture}}
\newcommand{\chaineinformation}{\begin{tikzpicture}[font=\small,>=Latex,node distance=5mm,bloc/.style={draw,minimum width=25mm,minimum height=11mm,fill=yellow!30,align=center}]
\node[bloc](a){\dots};\node[bloc,right=of a](b){\dots};\node[bloc,right=of b](c){\dots};\draw[->,very thick,SIblue](a)--(b);\draw[->,very thick,SIblue](b)--(c);\end{tikzpicture}}

\setlength{\parindent}{0pt}
\begin{document}
\InterroHeader{Loi entrée--sortie}{10}

\textbf{Objectif :} déterminer la loi entrée--sortie en position entre la translation du vérin et la rotation du levier 1.

\medskip
Le dessin ci-dessous représente la vue en coupe d'une pompe et son modèle cinématique. On pose
\[
\overrightarrow{OA}=R\,\vec x_1,\qquad \overrightarrow{AB}=x(t)\,\vec x_2,\qquad \overrightarrow{OB}=L\,\vec x.
\]
On note $\alpha=(\vec x,\vec x_1)$ et $\beta=(\vec x,\vec x_2)$.

\begin{center}
\begin{tikzpicture}[scale=1,>=Latex]
  \coordinate (O) at (0,0); \coordinate (B) at (8,0); \coordinate (A) at (2.3,1.45);
  \draw[->] (-.5,0)--(8.8,0) node[right]{$\vec x$};
  \draw[->] (O)--++(75:2.2) node[above]{$\vec y$};
  \draw[very thick,red] (O)--(A) node[midway,above left]{1};
  \draw[very thick,blue] (A)--($(A)!0.62!(B)$) node[midway,above]{2};
  \draw[very thick,brown] ($(A)!0.55!(B)$)--(B) node[midway,below]{3};
  \filldraw[fill=white,thick] (O) circle(2.2pt) node[below left]{$O$};
  \filldraw[fill=white,thick] (A) circle(2.2pt) node[above]{$A$};
  \filldraw[fill=white,thick] (B) circle(2.2pt) node[above]{$B$};
  \draw (O) ++(0:0.6) arc[start angle=0,end angle=32,radius=.6] node[pos=.7,right]{$\alpha$};
  \draw (B) ++(180:0.65) arc[start angle=180,end angle=166,radius=.65] node[pos=.55,left]{$\beta$};
  \draw[thick] (-.35,-.22)--(.35,-.22); \foreach \x in {-.3,-.15,...,.3}{\draw (\x,-.22)--++(-.12,-.18);} 
  \draw[thick] (7.65,-.22)--(8.35,-.22); \foreach \x in {7.7,7.85,...,8.3}{\draw (\x,-.22)--++(-.12,-.18);} 
\end{tikzpicture}
\end{center}

\question{Compléter le graphe de liaisons ci-dessous.}{2}
\begin{center}
\begin{tikzpicture}[>=Latex]
\foreach \n/\x/\y in {0/0/0,1/4/0,2/4/-2.2,3/0/-2.2}{\node[draw,circle,minimum size=8mm] (n\n) at (\x,\y) {\n};}
\draw (n0)--node[above]{Pivot d'axe $(O,\vec z)$}(n1);
\draw (n0)--node[left]{Pivot d'axe $(B,\vec z)$}(n3);
\draw (n3)--node[below]{\dotfill\hspace{28mm}}(n2);
\draw (n1)--node[right]{\dotfill\hspace{28mm}}(n2);
\end{tikzpicture}
\end{center}

\newpage
\InterroHeader{Loi entrée--sortie}{10}
\question{Donner les paramètres des mouvements d'entrée et de sortie.}{1}
\begin{reponse}[18mm]\end{reponse}
\question{Écrire, sous forme vectorielle, la fermeture géométrique. Ne pas résoudre.}{1}
\begin{reponse}[25mm]\end{reponse}
\question{Trouver une autre méthode permettant d'obtenir le résultat plus rapidement.}{1}
\begin{reponse}[30mm]\end{reponse}

\question{Donner la liaison équivalente des liaisons en parallèle suivantes.}{5}
\renewcommand{\arraystretch}{1.2}
\begin{tabular}{|p{.31\linewidth}|p{.31\linewidth}|p{.31\linewidth}|}\hline
\centering
\begin{tikzpicture}[scale=.55]\draw[thick](-2,0)--(2,0);\draw[very thick,blue](-1.5,0)--(1.5,0);\draw[red,thick](-1.5,0)circle(.25);\draw[red,thick](1.5,0)circle(.25);\node at (-1.5,.55){$A$};\node at (1.5,.55){$B$};\node at (0,.4){1};\node at (-2.2,.4){0};\end{tikzpicture}
&\centering
\begin{tikzpicture}[scale=.55]\draw[very thick,red](-1.8,-.6)--(1.8,.6);\draw[very thick,blue](-1.4,-.5)--(1.4,.5);\draw[thick](-1.2,-.4)circle(.22);\draw[thick](1.2,.4)circle(.22);\end{tikzpicture}
&\centering
\begin{tikzpicture}[scale=.55]\draw[very thick,blue](-1.7,.7)--(1.7,-.7);\draw[very thick,green!60!black](-1.4,-.6)--(1.4,.6);\draw[thick](-1.2,-.5)circle(.22);\draw[thick](1.2,.5)circle(.22);\node at (-1.55,-.85){$A$};\node at (1.55,.85){$B$};\end{tikzpicture}\\[15mm]\hline
\centering Liaison équivalente : &\centering Liaison équivalente : &\centering Liaison équivalente :\\[16mm]\hline
\centering
\begin{tikzpicture}[scale=.55]\draw[very thick,blue](-1.5,0)--(1.5,0);\draw[red,thick](-1.5,0)circle(.25);\draw[red,thick](0,0)circle(.25);\draw[thick](1.5,-.5)--(1.5,.5);\node at (-1.5,.55){$O$};\node at (0,.55){$A$};\end{tikzpicture}
&\centering
\begin{tikzpicture}[scale=.55]\draw[thick](-1.8,0)rectangle(1.8,.55);\draw[thick](-1.5,.28)circle(.22);\draw[thick](1.5,.28)circle(.22);\node at (-1.5,.8){$O$};\node at (1.5,.8){$P$};\end{tikzpicture}
&\\[15mm]\hline
\centering Liaison équivalente : &\centering Liaison équivalente : &\\[22mm]\hline
\end{tabular}

\end{document}
